


\section{Special rank one matrices in Vec($S$)}  %IV


In this section we will focus on the vector spaces generated by either the matrix semigroup or its kernel. 
Any rank-one matrix with identical rows in vec($S$) will be referred to as rank-one special, or ROSP.  
A ROSP in vec($K$) will be a KROSP.  
Note, that the ROSP’s form a vector space of dimension at most $n$.  
In general, we will identify a ROSP with its first row.  
Observe that if $M$ is any ROSP, then $MN$ is a ROSP for any matrix $N$ in vec($S$).
In addition, if $w$ is any word in $S$, $wM=M$.
That is, $M$ is a right zero for elements in $S$.


A \textit{range vector} is a vector consisting of 0's and 1's with 1's corresponding to the range of an idempotent $e$ in $K$, that is, a vector whose components indicate the non-zero columns of $e$. 
A range matrix is a KROSP whose rows equal some range vector.  
If there is no chance of confusion, for simplicity, we identify the range vector with its corresponding KROSP.


\begin{proposition}  \label{prop:5.2}
The vector space spanned by the KROSPs equals: 
\begin{enumerate}
\item[{\normalfont 1.}]  Span of $\{\Om w: w \in S \}$.
\item[{\normalfont 2.}]  Span of  $\{Mk: k \in K\}$, where $M$ is any ROSP in vec($S$) satisfying  $M\bu \not= 0$.
\item[{\normalfont 3.}]  Span of the range matrices.
\end{enumerate}
\end{proposition}


\begin{proof}    For \#1, first note that $\Om$ is an element of vec($K$), and therefore a KROSP (see the proof of Theorem [7] in \cite{BM}).  Since $K$ is an ideal, $kw \in K$, for any  $w \in S$.  
So $\Om w \in$ vec($K$) and hence a KROSP. 
Equality will follow from \#2.


For \#2, for  $k \in K$, $Mk \in$ vec($K$)  and so is a KROSP.   
If the set  $\{Mk : k \in K\}$ does not span the KROSPs, then there is a KROSP, $X$, with first row $\bx$ such that 
\[
  Mk\bx^t=0, \ \hbox{ for all } \ k \in K.
\]


Let  $M\bu = c\bu$, with  $c \not= 0$.  
Since $X$ itself is a KROSP, $cXX^t = MXX^t = 0$.  
So $X = 0$ and the result follows. 


For \#3, by the proof of Theorem 7 in \cite{BM}, $\pi k$ is a multiple of a range vector for every $k \in K$.  
So the result follows from \#2, taking  $M=\Om$. 
\qed
\end{proof}


The constant functions play a special role in any investigation of the road coloring problem. 
Let $C_i$ denote the matrix of the constant function $i$; that is for $C_i$, the $i$th  column is all ones.  
If $S$ contains the constants  $C_i$, they must be in $K$ due to minimality of the rank of the kernel.
Clearly, any element in the linear span of the constants  $C_i$, $i = i,\ldots,n$ is a KROSP.   
So the KROSPs have dimension $n$ if and only if vec($K$) contains the constants  $C_i$.  
In fact, we can say more.


\begin{proposition}  	\label{prop:5.3}
The KROSPs have dimension n if and only if the rank of the kernel is one.
\end{proposition}


\begin{proof}
One direction is immediate.  
For the other, assume that the KROSPs have dimension $n$. 
By the above, vec$(K)$ contains the constants $C_i$.


Now  $C_i = \sum\limits_{k \in K} \alpha_k k$.
Choose an idempotent $e \in K$.  
Then $eC_i e = \sum\limits_{k=K} \alpha_k eke$.  
Since $K$ is a completely simple semigroup (see \cite{CP}), then $eC_i e= C_j$  is in vec$(G)$ where $G = eKe$ is the subgroup for which $e$ is the local identity element (again \cite{CP}), and $ie=j$.


Suppose the rank of $K$ is $r$.  
Then $G$ is a group of block permutation matrices of rank $r$.  
In particular, the matrices of $G$ have $r$ nonzero columns.  
Notice $\Om C_j = C_j$.
On the other hand, $\Om C_j = \sum\limits_{k=K} \alpha_k \Om eke = C_j$.
But, as stated above, the proof of Theorem 7 in \cite{BM} shows that, for each $k \in K$, 
$\Om eke$ is a fixed ROSP whose rows are multiples of the range vector for the group $G$.  
Thus $C_j = \Om eke \sum\limits_{k=K} \alpha_k$ where $\Om eke$ is a ROSP with $r$ non-zero rows.  
But this is only possible if $r =1$. 
\qed
\end{proof}
